\chapter{数学运算能力的概述}


% \section{数学运算能力的概念}

\section{数学运算}

% 《普通高中数学课程标准（2017年版）》中明确指出数学运算是数学学科核心素养之一。数学运算是指在明晰运算对象的基础上，依据运算法则解决数学问题的素养。主要包括：理解运算对象，掌握运算法则，探究运算思路，选择运算方法，设计运算程序，求得运算结果等。数学运算主要表现为：理解运算对象，掌握运算法则，探究运算思路，求得运算结果。通过高中数学课程的学习，学生能进一步发展数学运算能力；有效借助运算方法解决实际问题；通过运算促进数学思维发展，形成规范化思考问题的品质，养成一丝不苟、严谨求实的科学精神\cite{教育部全日制义务教育数学课程标准}。

《普通高中数学课程标准（2017 年版）》中明确指出数学运算是数学学科核心素养之一\cite{教育部全日制义务教育数学课程标准}。数学运算是指在明晰运算对象的基础上，依据运算法则解决数学问题的素养。主要包括：理解运算对象，掌握运算法则，探究运算思路，选择运算方法，设计运算程序，求得运算结果等。数学运算主要表现为：理解运算对象，掌握运算法则，探究运算思路，求得运算结果\cite{教育部全日制义务教育数学课程标准}。在中学数学课上，可以使学生的数学计算能力得到进一步的发展；利用算术的高效运用求解问题；通过运算促进数学思维发展，形成规范化思考问题的品质，养成一丝不苟、严谨求实的科学精神\cite{教育部全日制义务教育数学课程标准}。

\section{数学运算能力}

数学运算能力就是数学学习时的计算能力，在数学教学中，学生的数学素养直接关系到他们的学业成绩。
数学计算不是单一的，它是将思维和计算技巧有机地融合在一起，而计算的能力则是通过对数学的学习来实现的。学生具备了一定的数学计算能力，主要体现在对问题进行快速的解析与求解，并针对不同的问题，选取适当的解题方式，以迅速地进行问题的分析与求解，并从中获取相关的知识与技巧。


% \section{培养西藏高中生数学运算能力的意义}


% 高中数学运算中有很多知识点存在，学生们接受知识的能力存在差异，在数学运算方面学生很容易出现问题，而数学运算能力的高低直接影响学生对数学学科的学习。首先，数学运算能力很大程度上影响学生学习数学的积极性和热情\cite{康恬雨西藏高中生数学运算能力的发展现状与培养策略研究}。运算能力差的学生经常做错题，甚至以前做过的题也会出现新的计算错误，削弱了学生学习数学的兴趣。其次，数学运算能力也影响着学生解决问题的速度。高考是限时考试，运算能力差的同学解题速度慢，容易出现考试时间不够用而做不完题的情况，最后的成绩自然差强人意。此外，数学运算能力还决定了课堂效率。运算能力差的同学上课容易跟不上教师的思路，而高中作业繁多，学生很难挤出时间课后补上；而整体运算能力差容易导致教师的上课效率低，课上讲不完，会拖堂或补课，又会耽误学生很多课下时间。

% 所谓“学好数理，行走江湖，无所畏惧”，显然，学习好了也是学习其它科目的必要条件，而逻辑严密的逻辑思考对于解决日常问题也有很大的好处。

% 培养学生的数学计算技能，有利于提高他们的数学思考，使他们的思想素质得到提高，使他们养成一丝不苟、严谨求实的科学态度\cite{龙奕羽素质教育视野中高数教学与课程思政的融合}。